\centerline{\bf \Huge Графики}


\begin{flushright}
        \scriptsize{\textbf{„Теперь мы в финале, мы должны победить, и не можем думать, что работа закончена. Мы отдадим все, и мы убеждены, что мы сможем сделать это."\\
         CR7}}
        
\end{flushright}
\textbf{План полного исследования функции и её построения:}\\
1) Нахождение области определения, по возможности, области значений и, если $\mbox{требуется}$, её обязательного дополнения после полного исследования.\\
2) Определение чётности, нечётности, общего вида.\\
3) Проверка наличия периодичности.\\
4) Нахождение нулей функции и интервалов знакопостоянства.\\
5) Нахождение экстремумов и промежутков возрастания, убывания.\\
6) Нахождение точек перегиба и промежутков выпуклости, вогнутости.\\
7) Доказательство непрерывности или определение точек разрыва.\\
8) Нахождение всех асимптот.\\
9) Дополнительные точки для точности построения.

\problem{\textbf{1}} $f(x)=\sqrt{x^{2}-\left|x\right|+3}$

\problem{\textbf{2}} $f(x)=\dfrac{(x-2)^{3}}{2(x-1)^{2}}$

\problem{\textbf{3}} $f(x)=\dfrac{\left|x\right|}{\sqrt{x^{2}+x}}$

\problem{\textbf{4}} $f(x)=\sqrt[3]{\left(x-1\right)^{2}(\left|x+1\right|)}$

\problem{\textbf{5}} $f(x)=\dfrac{e^{x-1}}{\left|x+1\right|}$

\problem{\textbf{6}} $f(x)=\ln\dfrac{x^{2}+2}{x^{2}-3}$

\problem{\textbf{7}} $f(x)=\left\{\begin{array}{l}
\dfrac{2}{x+7}+1, x<-4 \\
\dfrac{4}{3}\left|x-1\right|-5, -4< x<6\\
0,5x-\dfrac{1}{x-13}, 6\leq x 
\end{array}\right.$

\problem{\textbf{8}} $f(x)=(x-1)^{2}e^{-\dfrac{x^{2}}{5}}$
\centerline{\textbf{240 минут}}

